% =====================================================================
% Portada de la memoria (diseño de la plantilla de la UAX, hecha en TikZ).
% Redactada el 06/11/2026 (sesión R4).
% Datos completados el 08/11/2026. Si falta figuras/logo_uax.png sale un recuadro rojo.
% =====================================================================
\definecolor{uaxazul}{HTML}{2F5496}
\newcommand{\pendiente}[1]{\textcolor{red}{\textbf{[#1]}}}

% --- Datos de la portada ---
\newcommand{\portadaCurso}{MUIA 2026}
\newcommand{\portadaCentro}{Business Tech}
\newcommand{\portadaAutor}{Guillermo Soriano Pérez}
\newcommand{\portadaDirector}{Ángel Manuel Rayo Acevedo}
\newcommand{\portadaFecha}{Noviembre de 2026}

\begin{titlepage}
\thispagestyle{empty}
\begin{tikzpicture}[remember picture,overlay]
	% marco azul
	\draw[uaxazul,line width=2.2pt]
		([xshift=1.4cm,yshift=-1.4cm]current page.north west) rectangle
		([xshift=-1.4cm,yshift=1.4cm]current page.south east);
	% franja lateral
	\fill[uaxazul]
		([xshift=1.4cm,yshift=-1.4cm]current page.north west) rectangle
		([xshift=4.6cm,yshift=1.4cm]current page.south west);
	\node[white,font=\bfseries\Huge] at ([xshift=3.0cm,yshift=-3.0cm]current page.north west) {UAX};
	\node[white,font=\bfseries\large,rotate=90] at ([xshift=3.0cm]current page.west) {TRABAJO FIN DE MÁSTER};
	\node[white,font=\bfseries\normalsize,align=center] at ([xshift=3.0cm,yshift=3.0cm]current page.south west) {\portadaCurso};
\end{tikzpicture}

% contenido, desplazado a la derecha de la franja
\hspace*{3.4cm}\begin{minipage}{0.64\textwidth}
	\raggedright
	\vspace*{0.3cm}
	\IfFileExists{figuras/logo_uax.png}{%
		\includegraphics[width=5.5cm]{figuras/logo_uax.png}}{%
		\fbox{\pendiente{logo de la UAX: figuras/logo\_uax.png}}}

	\vspace{1.2cm}
	{\color{uaxazul}\large\bfseries Universidad Alfonso X el Sabio}\\[0.2cm]
	{\color{uaxazul}\large \portadaCentro}\\[0.2cm]
	{\color{uaxazul}\large Máster Universitario en Inteligencia Artificial}

	\vspace{2.4cm}
	{\color{uaxazul}\Large\bfseries Previsión a corto plazo de la demanda eléctrica
	peninsular: comparación de modelos clásicos y de aprendizaje automático con
	evaluación de intervalos de predicción}

	\vspace{2.4cm}
	{\large\textbf{Autor:} \portadaAutor}\\[0.3cm]
	{\large\textbf{Director:} \portadaDirector}\\[0.3cm]
	{\large \portadaFecha}
\end{minipage}
\end{titlepage}
